\section{应用价值}

本章讨论系统产品化的\textbf{目标价值与情景测算};数据来源、字段质量与分析可信边界详见《SciScope 数据分析报告》。真实效率、经济与社会效益仍需合作单位试点和前后对照验证。

\subsection{社会效益}
\begin{itemize}
\item \textbf{科研提效与公平}:把分散在多源、多语言的文献统一为可问答知识底座,目标是降低文献调研工具的使用门槛;是否形成可量化的信息可及性改善,仍需真实用户试点验证。
\item \textbf{可信与可审计}:关键可验证回答均附引用论文(evidence-grounded);\textbf{论断核查}输出支持/反驳/证据不足并保留拒答边界,\textbf{工具校验门}在调用前拦截模型编造的论文主键——从生成与调用两端共同抑制大模型幻觉与捏造,支撑\textbf{科学决策}与严肃的立项查新、综述撰写,减少误导性结论传播。该结论是工程能力边界,不等于临床或监管级判断。
\item \textbf{知识传承与脉络梳理}:主题演化、关键词趋势与合作网络帮助新人快速理解领域发展脉络,促进知识传承。
\item \textbf{本地部署边界}:数据、索引和可选本地大模型可留在本地,减少查询正文外发;涉密或敏感场景上线前仍须完成单位级安全、权限与合规评审,本报告不宣称已取得认证。
\item \textbf{开放与可复用}:工具经 MCP 协议开放,接地语料检索与论断核查可被兼容客户端复用,具备减少重复接口开发的潜力;实际复用规模仍待外部接入验证。
\end{itemize}

\begin{tableblock}
\begin{tabularx}{0.96\textwidth}{>{\bfseries}p{28mm}X X}
\toprule
受益群体 & 具体行动 & 可验证支撑 \\
\midrule
青年研究者 / 研究生 & 用中文问题直达英文前沿,快速获得代表论文、趋势和引用证据 & 固定中文主题样例上的跨语言检索已验证; 回答附证据卡与支持/反驳/证据不足边界 \\
科研管理者 & 基于热点、生命周期、合作网络制定学科布局和指南方向 & 数据分析报告中的关键词演化、作者社区和质量边界 \\
高校图书馆 / 情报部门 & 将年度文献巡检沉淀为可复现流程和可复用报告模板 & \texttt{make} 流水线、报告 PDF、评测证据包 \\
企业研发 / 技术情报团队 & 用趋势、推荐和图谱筛选技术路线、竞品方向与潜在合作对象 & 四路自动多基线表明 `current` 相对 `dense` 在同领域率、共享关键词、多样性和新颖度代理上更高,但语义相似度略低且延迟更高; 真实用户价值仍待验证 \\
涉密或内网机构 & 在本地部署语料、索引、嵌入器和生成模型,减少敏感查询外发 & PostgreSQL + pgvector + 本地 7B / 可替换生成层 \\
\bottomrule
\end{tabularx}
\end{tableblock}

\subsection{经济效益}
\begin{itemize}
\item \textbf{效率目标}:系统把逐篇检索、初筛和证据整理组织成可复现流程;实际节省工时和错误降低幅度尚未实测,须由用户任务前后测确认。
\item \textbf{成本结构}:核心能力沉淀在本地索引与模型文件;本地 7B 可减少云端 API token 可变费用,但仍产生硬件折旧、电力、运维与人员成本,不能写成总边际成本为零。
\item \textbf{决策支持潜力}:描述性趋势、热点识别和争议证据链可作为研发立项与技术路线讨论的输入;是否降低试错成本仍待真实项目验证,且当前趋势不得写成预测能力。
\item \textbf{合作拓展}:作者合作网络辅助识别潜在合作者与团队,促进产学研对接与资源整合。
\item \textbf{扩展低成本}:新增能力靠\textbf{插入 MCP 服务}或派生专员子智能体获得,而非重新开发;能力边界增量扩展,边际开发成本低。
\item \textbf{数据要素服务出口}:MCP 与 OpenCode 顺序调用已证明 SciScope 不只是一个终端助手,还可以作为被其他智能体调用的科研证据服务出口。
\end{itemize}

\subsubsection{经济效益量化(情景测算)}
为便于评委直观判读,下面给出\textbf{情景测算}(基于公开订阅定价与典型科研流程工时假设),用于说明本地化方案的成本结构差异。

\begin{tableblock}
\begin{tabularx}{0.96\textwidth}{>{\bfseries}p{30mm}X X}
\toprule
方案 & 年度可变成本(以 100 名研究员估) & 数据与自主可控 \\
\midrule
商用文献 AI 订阅 & 按 \$15/月/人估,约 \textbf{12--18 万元/年}, 随人数线性增长 & 查询与文献记录依赖第三方平台,数据主权受平台条款约束 \\
通用大模型 API 问答 & 按重度文献问答用量估,\textbf{数万至十余万元/年}, 随调用量增长 & 查询内容需发送至外部生成服务 \\
\rowcolor{SciPanel} SciScope 本地部署 & 建库与索引可复现;本地推理无云端 token 费,但保留硬件、电力、运维与人力成本 & 数据、索引和可选生成模型可本地部署;正式安全结论待单位评审 \\
\bottomrule
\end{tabularx}
\end{tableblock}

\begin{tableblock}
\begin{tabularx}{0.96\textwidth}{>{\bfseries}p{30mm}X X}
\toprule
科研环节 & 传统人工 & SciScope 辅助(情景估计) \\
\midrule
立项查新 / 综述前期调研 & 单次约 40--80 工时(逐篇检索、去重、筛读) & 语义问答 + 证据直达,估\textbf{压缩 30\%--50\%} 前期工时 \\
热点 / 趋势研判 & 数天人工巡检多个数据库(情景假设) & 脚本可生成描述性榜单;实际任务耗时待试点计时 \\
相似论文 / 合作者发现 & 人工滚雪球式查找 & 推荐与图谱即时给出候选 + 可解释因子 \\
\bottomrule
\end{tabularx}
\end{tableblock}

\textbf{单次问答成本(基于 token 遥测的核算口径)}:本系统在每轮智能体回答中返回估算 token 用量与终止原因;该数值来自运行时对消息内容的确定性估算,用于\textbf{逐次记录、复盘和成本核算},而不是本报告中虚构一组固定成本。若接入云端 DeepSeek 等 OpenAI 兼容模型,组织可把 token 用量乘以当期公开 API 单价进行\textbf{情景估算};若使用本地 7B 推理,则主要成本转为本机算力、电力和运维投入。本项目已经具备成本可观测接口,但尚未开展正式用户规模化成本实验。

\begin{plainbox}
\textbf{投入产出 / 回本口径}:上述数值按公开订阅区间与典型流程假设做情景估算,用于说明成本结构变化方向,不代表实测收益。当前只能确认本地化流程具备标准化检索、初筛和证据整理的工程条件;效率和回本周期须由真实业务计时与完整运维成本核算验证。
\end{plainbox}

\begin{plainbox}
\textbf{测算边界(据实声明)}:以上为\textbf{情景测算},依据公开订阅定价区间与典型科研流程工时假设推得,用于说明成本结构与效率方向;\textbf{具体金额、工时与百分比须结合目标机构的真实业务计时或用户实验核对},本报告不声称已实测,亦不夸大未验证的经济数字。
\end{plainbox}

\subsection{推广前景}
\begin{itemize}
\item \textbf{高校/科研院所}:学科情报中心、研究生文献调研、综述与查新服务。
\item \textbf{企业研发}:技术情报监测、竞品与专利分析、研发选题。
\item \textbf{科研管理}:基金评审辅助、领域态势研判、合作网络分析。
\item \textbf{可扩展性}:架构与领域无关,更换语料即可迁移到任意学科或企业内部文档库;模型可替换、流水线可复现,便于产品化与私有化部署。
\end{itemize}

\begin{reportbox}{差异化价值小结}
相比"直接问通用大模型",SciScope 的不可替代之处在于:\textbf{基于真实语料、有出处可追溯、本地可复现、领域可迁移、大模型可替换}。
它不是又一个聊天机器人,而是一套\textbf{把机构自有文献资产转化为可信科研生产力}的智能体系统。
\end{reportbox}
